\section{Chapitre~\ref{sec:pain}. La douleur}

Les capteurs de la douleur, les deux fils, le portillon de la moelle épinière, le bouton de volume descendant, la sensibilisation et la capture de l'attention sont mesurés directement; les formules du portillon et de la sensibilisation sont des modèles simplifiés du livre; la règle selon laquelle la machine choisit le volume de l'alarme est une hypothèse.

{\footnotesize
\setlength{\tabcolsep}{3pt}\renewcommand{\arraystretch}{1.15}
\begin{longtable}{>{\raggedright}p{0.5cm}>{\raggedright}p{4.0cm}>{\raggedright}p{5.6cm}>{\raggedright}p{2.3cm}>{\raggedright\arraybackslash}p{1.7cm}}
\toprule
N\textsuperscript{o} & Affirmation & Expérience et ce qui est mesuré & Source & Statut \\
\midrule\endfirsthead
\toprule
N\textsuperscript{o} & Affirmation & Expérience et ce qui est mesuré & Source & Statut \\
\midrule\endhead
1 & Seuil élevé: les seuils d'échauffement des différents capteurs de la douleur vont de $41^\circ$ à $46$--$48^\circ$ & Enregistrement de fibres isolées. Peau de singe: seuil médian de $41^\circ$ pour les capteurs lents (C), $46^\circ$ pour les capteurs rapides de type II, plus de $53^\circ$ pour le type I. Peau humaine: $40.7^\circ$ pour les capteurs C qui répondent aussi à la pression, $48^\circ$ pour ceux qui n'y répondent pas. & \cite{Treede1995heat, Weidner1999cfib} & mesuré \\
2 & Les capteurs de la douleur s'adaptent bien moins que les capteurs du toucher et deviennent plus sensibles après une lésion légère; un affaiblissement après un fort échauffement est mesuré pour un seul type & Brûlure légère de la peau ($50^\circ$, $100\,\text{s}$): après $5$--$10$~min, le seuil de douleur chez l'homme et le seuil des capteurs C chez le singe sont abaissés de $2$--$6^\circ$; les autres capteurs cutanés ne présentent pas ce décalage. Chez les capteurs rapides de type II, la réponse est supprimée après un fort échauffement. La comparaison de l'accoutumance avec les capteurs du toucher est une estimation générale, non étayée ici par une expérience propre. & \cite{LaMotte1982hyper, Treede1995heat} & mesuré \\
3 & Deux fils: rapide, $5$--$30$~m/s; lent, $0.5$--$2$~m/s; temps d'arrivée $t=L/v$~\eqref{eq:paintime} & Vitesse de conduction: capteurs rapides de la douleur du singe, en moyenne $15$ et $25$~m/s (deux types); capteurs C humains, $0.8$--$1.0$~m/s. Pour $L=1$~m, cela fait quelques dizaines de millisecondes et environ une seconde. & \cite{Treede1995heat, Weidner1999cfib} & mesuré \\
4 & La douleur arrive deux fois: d'abord rapide et précise, puis sourde et durable & Temps de réaction à l'échauffement de l'avant-bras humain: de $1100$ à $400\ms$ quand la température monte; un fort échauffement de la peau velue donne une double douleur, celui de la paume non. Seules des fibres plus rapides que $6$~m/s peuvent porter la première douleur; par la latence, elle correspond aux capteurs rapides de type II, absents de la paume. Le partage des rôles (\og où\fg{} et \og à quel point c'est grave\fg{}) est une interprétation. & \cite{CampbellLaMotte1983first, Treede1995heat} & mesuré \\
5 & Portillon: le neurone de sortie de la moelle reçoit la douleur et le toucher, l'intermédiaire inhibiteur est excité par le toucher (fig.~\ref{fig:gate}) & Le schéma du portillon a été proposé comme théorie. Souris, marquage génétique et suppression de groupes de neurones: les neurones excitateurs à somatostatine sont nécessaires à la douleur mécanique; les neurones inhibiteurs à dynorphine empêchent l'entrée des capteurs du toucher de les atteindre; sans ces neurones, un toucher léger provoque de la douleur. & \cite{MelzackWall1965, Duan2014} & mesuré \\
6 & Le toucher atténue la douleur & Homme, impulsions laser douloureuses sur la main: un toucher simultané réduit la douleur ressentie et la capacité à distinguer l'énergie des impulsions; l'effet faiblit avec la distance entre le point de toucher et le point de douleur au sein d'un même segment cutané. & \cite{Mancini2014touch} & mesuré \\
7 & Hypothèse de la théorie du portillon: une douleur forte éteint elle-même l'intermédiaire inhibiteur, et le portillon s'ouvre grand & Présenté dans le chapitre comme une hypothèse de la théorie originale. Le travail sur la souris décrit comment le portillon empêche le toucher d'entrer dans le canal de la douleur; l'extinction de l'intermédiaire par la douleur n'y est pas montrée. & \cite{MelzackWall1965} & hypothèse \\
8 & Portillon~\eqref{eq:gate}: seuil de $0.18$ sans toucher, $0.86$ avec toucher, $0.11$ pour $g=2$; pour $\nu=1$, le toucher réduit la sortie de $1$ à $0.25$, pour $\nu=2$ il est sans effet; le toucher seul ne passe pas (fig.~\ref{fig:gatecurves}) & Le calcul avec \texttt{models/\allowbreak pain\_\allowbreak gate.py} et les poids du texte reproduit tous ces nombres. & \texttt{models/\allowbreak pain\_\allowbreak gate.py} & modèle \\
9 & Les voies descendantes du tronc cérébral modifient le gain du portillon dans les deux sens & Rat, bulbe rachidien, enregistrement pendant l'échauffement de la queue: certains neurones déchargent avant le retrait (cellules \og ON\fg{}), d'autres se taisent (cellules \og OFF\fg{}); une faible stimulation de cette région supprime le réflexe. Revue: les mêmes circuits facilitent et inhibent la transmission de la douleur, et les opioïdes agissent par eux. & \cite{Fields1983rvm, Fields2004} & mesuré \\
10 & L'attente d'un soulagement atténue la douleur par le canal opioïde & Patients souffrant d'une douleur aiguë: chez ceux dont l'attente d'un soulagement avait réduit la douleur, le blocage des récepteurs opioïdes a ramené la douleur au niveau des autres; chez les autres, le blocage ne changeait pas la douleur. & \cite{Levine1978} & mesuré \\
11 & Le cerveau choisit le volume de l'alarme selon l'intérêt de l'interruption & Il n'existe pas d'expérience mesurant la règle du choix du volume. & --- & hypothèse \\
12 & Sensibilisation: après une lésion, la sortie du portillon augmente et le seuil baisse, en partie grâce à la moelle épinière; elle persiste après l'arrêt du stimulus nocif & Rat: après une lésion cutanée, le seuil du réflexe de flexion baisse et la réponse augmente; l'analyse électrophysiologique montre qu'une partie de cette hausse naît dans la moelle elle-même. Un stimulus nocif provoque des troubles durables de la sensibilité qui survivent au stimulus. Le chapitre ne donne pas le temps de décroissance exact. & \cite{Woolf1983} & mesuré \\
13 & La sensibilisation~\eqref{eq:sensitization} est une réaction positive; si la boucle est forte, l'alarme peut se verrouiller après la guérison & Découle du modèle~\eqref{eq:sensitization} (exercice en fin de chapitre). Aucune expérience ne montre qu'une douleur persistante après guérison soit une boucle verrouillée de ce type précis. & déduction du chapitre & hypothèse \\
14 & La douleur est une récompense négative, et l'apprentissage par la douleur suit l'erreur de différence temporelle & IRM, sujets humains, chaînes de signes annonciateurs de la douleur: l'activité du striatum ventral et de l'insula antérieure coïncide avec les signaux d'apprentissage par différence temporelle. Singe: une partie des neurones \nt{DOPA} est excitée par un souffle d'air désagréable et son signe annonciateur, une autre partie est inhibée. & \cite{Seymour2004, Matsumoto2009} & mesuré \\
15 & La douleur interrompt le travail en cours & Sujets humains, stimulus électrique douloureux et sondes sonores: la réponse à une sonde présentée $100\ms$ après le début de la douleur est nettement ralentie; à $1.5\,\text{s}$, il n'y a plus de différence avec le contrôle. Une revue rassemble ces expériences dans un modèle de capture de l'attention. & \cite{Crombez1996disrupt, EcclestonCrombez1999} & mesuré \\
16 & Le retrait se referme dans la moelle épinière & Rat: les neurones des couches profondes de la corne dorsale reproduisent le motif spatial du réflexe de retrait de muscles isolés; on ne parvient pas à les exciter de façon antidromique depuis la région cervicale: ce sont donc des intermédiaires spinaux. & \cite{Schouenborg1995wr} & mesuré \\
\bottomrule
\end{longtable}}
